

\begin{proposition}[D.10 --- Proved: the literal least quotient on locally soluble curves]
\label{D-literal-minimum}
For a fixed $C/K$ with $C(K_v)\ne\varnothing$ at
every place, the least quotient of $\pi_1(C)$
over $G_K$ for which every global section is
locally point-induced is $G_K$ itself.
Its geometric kernel is trivial.
\end{proposition}
\begin{proof}
The quotient is the structure map
$\pi_1(C)\twoheadrightarrow G_K$.
Its structure map to $G_K$ is the identity,
whose only section is the identity. After
pullback to $G_{K_v}$ the structure map is
again the identity. For every $x_v\in C(K_v)$
the composite
\[
 G_{K_v}\xrightarrow{s_{x_v}}\pi_1(C_{K_v})
 \longrightarrow G_{K_v}
\]
is exactly that identity. Thus the unique
global section is locally point-induced.
Every quotient over $G_K$ maps to $G_K$,
which proves minimality in the factorization
order. Conversely, at a place with no
local point the image of the point map
is empty, so this quotient does not have
the property.

The literal assertion consequently has
this answer on every locally soluble
curve, including every curve with a
rational point. The comparison sends
all its point-section classes to the
same class. Requiring recovery of
the original point-section classes
would impose an additional condition
on this very map; we have not inserted
such a condition into the question.
A uniform selection for all curves,
including the locally insoluble ones,
is the remaining interpretation
recorded in Proposition~\ref{D-open}.
\end{proof}


\begin{theorem}[D.11 --- Refuted: a least quotient need not exist]
\label{D-no-least-quotient}
For a fixed proper curve $C/\mathbb Q$, order arithmetic
quotients by their specified continuous epimorphisms
$E_C\twoheadrightarrow Q$ over $G_{\mathbb Q}$:
a smaller quotient is obtained by a commuting further
epimorphism. Say that $Q$ has property $\mathcal P_C$
when every global section of $Q\to G_{\mathbb Q}$ is,
at every place $v$, in the image of the actual map
$C(\mathbb Q_v)\to\operatorname{Sect}(Q_v)/\ker(Q_v\to G_v)$.

There is no least quotient with $\mathcal P_C$ for
the explicit smooth projective genus-five curve
\[
 C:\quad y^2=-(u^6+1)(u^6+4).
 \tag{D11.1}
\]
There are two quotients $Q_1,Q_2$, each with geometric
kernel $C_2$, which have no global sections and hence
have $\mathcal P_C$. Their only common further quotient
over $G_{\mathbb Q}$, with the indicated maps from $E_C$,
is $G_{\mathbb Q}$; that quotient fails $\mathcal P_C$.
Together with Proposition~\ref{D-literal-minimum}, this
decides the least-quotient request in D(c) under the
literal stated property. The genus-five example also
has $C(\mathbb Q_2)=\varnothing$, so the failure persists
when only finite places are required.
\end{theorem}
\begin{proof}
We first construct the quotient maps, including their
arithmetic descent. Consider the smooth projective
genus-one curve
\[
 T:\quad y^2=-(x^2+1)(x^2+4).
\]
It has no real point: the right side is negative for
finite real $x$, and the leading coefficient $-1$
excludes real points at infinity. Its Jacobian has
rational full two-torsion. Indeed the four branch
points of $T\to\mathbb P^1$ are
$i,-i,2i,-2i$. The three nonzero geometric two-torsion
classes correspond to their three unordered partitions
into two pairs. The only nontrivial Galois permutation
is $(i,-i)(2i,-2i)$, which fixes every partition.
For this identification one may choose a branch
point as origin over $\overline{\mathbb Q}$:
the differences of the other branch points give
exactly the three nonzero points of order two.
Changing the origin translates all four labels and
preserves the three pair partitions. This also proves
the asserted Galois equivariance.

Consequently $G_{\mathbb Q}$ acts trivially on
\[
 \Delta_T/2\Delta_T
   \simeq H_1(T_{\mathbb C},\mathbb F_2)
   \simeq(\mathbb F_2)^2 .
\]
The comparison here is finite-cover comparison, with
the chosen complex embedding and geometric base point.
The action on the abelian quotient is independent of
the base-point path.

Complex conjugation acts freely and reverses
orientation on the torus $T(\mathbb C)$. Its quotient
is a Klein bottle: it is a closed nonorientable surface
with Euler characteristic zero, and its orientation
double cover is the torus. Thus its topological
arithmetic extension has presentation
\[
 E_{\mathbb R}^{\rm top}
   =\langle a,b\mid aba^{-1}=b^{-1}\rangle,\qquad
 \Delta_T^{\rm top}=\langle a^2,b\rangle\simeq\mathbb Z^2,
 \qquad a\longmapsto c,\ b\longmapsto1.
 \tag{D11.2}
\]
Completion identifies this extension with the etale
extension over $\mathbb R$. Define the two continuous
characters of the geometric group by
\[
\begin{array}{c|cc}
 &a^2&b\\ \hline
 \chi_1&1&0\\
 \chi_2&1&1
\end{array}
 \qquad\text{with values in }\mathbb F_2 .
 \tag{D11.3}
\]
They are independent and are invariant under all of
$G_{\mathbb Q}$ by the preceding two-torsion calculation.
Their kernels are therefore closed normal subgroups
of $E_T$. The quotients
\[
 E_T/\ker\chi_j\longrightarrow G_{\mathbb Q}
 \tag{D11.4}
\]
have kernel $C_2$ and trivial action on that kernel.
After restriction to $G_{\mathbb R}=C_2$, each is
$C_4\to C_2$: the lift $a$ of conjugation has square
equal to the nonidentity geometric element.
Every lift of conjugation has order four, so neither
restricted extension has a section. Restriction of a
global section would give such a section; hence
neither extension in (D11.4) has a global section.

The map
\[
 q:C\longrightarrow T,\qquad (u,y)\longmapsto(u^3,y)
 \tag{D11.5}
\]
extends to the smooth projective curves and has degree
three. The degree-twelve polynomial in (D11.1) is
squarefree, so $C$ is smooth of genus five. Pullback
on geometric $H^1(-,\mathbb F_2)$ is injective: the
trace map composed with pullback is multiplication
by $\deg q=3$, hence the identity over $\mathbb F_2$.
This trace identity follows as well from topological
Poincare duality and the degree of the continuous map
of compact oriented surfaces.

Pull back $\chi_1,\chi_2$ along the arithmetic
fundamental-group map induced by (D11.5), and denote
their independent geometric characters by
$\chi_1^C,\chi_2^C$. Put
\[
 Q_j=E_C/\ker\chi_j^C .
 \tag{D11.6}
\]
The map $q_*$ induces an isomorphism from each $Q_j$
to (D11.4): it is the identity on $G_{\mathbb Q}$ and
an isomorphism on their geometric kernels of order
two. An injective and surjective map of these
extensions follows directly by lifting an arithmetic
element and correcting it in the kernel. Thus the
two $Q_j$ have no global sections, as asserted.

Let $E_C\twoheadrightarrow Q$ be a common further
quotient through both maps in (D11.6). Its geometric
kernel is a quotient of each $C_2$. A nontrivial
such kernel would make both maps $C_2\to\ker(Q\to
G_{\mathbb Q})$ isomorphisms. The two composite
geometric characters would then have the same
kernel. This contradicts the independence in
(D11.3) and the injective pullback in (D11.5).
The geometric kernel of $Q$ is therefore trivial,
and the indicated map identifies $Q$ with
$G_{\mathbb Q}$. This quotient has its identity
section. There is no point of $C(\mathbb R)$ that
can induce its real localization, because
$C(\mathbb R)$ maps to the empty set $T(\mathbb R)$.
It fails $\mathcal P_C$, proving the nonexistence
of a least quotient.

For the additional finite-place assertion we verify
$T(\mathbb Q_2)=\varnothing$. For $x$ a two-adic unit,
$-(x^2+1)(x^2+4)$ has valuation one. For
$v_2(x)=1$ it has valuation three. For
$x=4z$ its quotient by four is the odd unit
\[
 -(16z^2+1)(4z^2+1)\equiv3\text{ or }7\pmod8 .
\]
For $v_2(x)<0$, remove the square $x^4$; writing
$z=x^{-1}\in2\mathbb Z_2$, the remaining unit is
\[
 -(1+z^2)(1+4z^2)\equiv3\text{ or }7\pmod8 .
\]
These cases also include $x=0$ in the third formula.
An odd two-adic unit is a square precisely when it is
one modulo eight. The points at infinity again
require a square root of $-1$. This proves the
claim, and (D11.5) gives $C(\mathbb Q_2)=\varnothing$.

The exact program checks squarefreeness, all residue
cases, the three branch partitions, and the finite
quotient of (D11.2)
\[
 C_4\times C_2,\qquad
 (a,b)\longmapsto a,\quad (a,b)\longmapsto a+2b
 \quad\text{in }C_4.
\]
Their two kernels generate the orientation kernel.
This is the finite certificate of the common-quotient
calculation; the global maps are (D11.4)--(D11.6).
The passage being answered is \S8, folio 33, scan 65,
with the explicit quotient follow-up in D(c).
\end{proof}


